\chapter{Statistical Learning Theory}
\label{cap:slt}

\begin{inbreve}
La Statistical Learning Theory (SLT) dà una risposta formale alla domanda centrale del corso: quando l'errore
sui dati di training è una buona stima dell'errore sui dati futuri? Il capitolo definisce rischio e rischio
empirico. Misura poi la complessità di uno spazio delle ipotesi con la VC-dimension. Infine presenta il bound sul
rischio e il principio della Structural Risk Minimization, da cui nascono le SVM.
\end{inbreve}

\section{Rischio e rischio empirico}
\label{sec:rischio}

Si vuole approssimare una funzione sconosciuta $f(\x)$, con target $d = f(\x)+\text{rumore}$. Idealmente si
vorrebbe minimizzare la \termine{funzione di rischio}, l'errore vero calcolato su \emph{tutti} i dati possibili:
\begin{equation}\label{eq:rischio}
R = \int L\big(d, h(\x)\big)\,dP(\x,d),
\end{equation}
dove $P(\x,d)$ è la distribuzione (sconosciuta) dei dati e $L$ una loss, ad esempio
$L(h(\x),d) = (d-h(\x))^2$. Si cerca $h\in H$ che minimizzi $R$. Si dispone però solo di un dataset finito
$TR = \{(\x_p,d_p)\}_{p=1}^{l}$. Si minimizza quindi il \termine{rischio empirico} (errore di training):
\begin{equation}\label{eq:rischio-empirico}
\Remp(h,TR) = \frac1l\sum_{p=1}^{l}\big(d_p - h(\x_p)\big)^2 .
\end{equation}
È il principio induttivo di \termine{Empirical Risk Minimization} (ERM). La domanda cruciale è:
\textbf{si può usare $\Remp$ per approssimare $R$?} Per rispondere serve una misura della complessità di $H$.

Le alternative per stimare l'errore atteso sono di due tipi:
\begin{itemize}
  \item \textbf{analitiche}: criteri AIC e BIC (\ingl{Akaike/Bayesian Information Criterion}, limitati ai
        modelli lineari nei parametri), MDL (\ingl{Minimum Description Length}), \textbf{SRM con la
        VC-dimension} (questo capitolo);
  \item \textbf{per ricampionamento}: cross-validation e bootstrap (capitolo~\ref{cap:validazione}).
\end{itemize}

\section{La VC-dimension}
\label{sec:vc}

Per un bound analitico serve misurare la \textbf{complessità} di $H$ (capacità, potere espressivo,
flessibilità). Va misurata anche quando $H$ è \textbf{infinito}, come nei modelli lineari con pesi reali: non si
può usare $\abs{H}$. L'idea di Vapnik e Chervonenkis è misurare \textbf{il numero di istanze distinte che $H$
riesce a discriminare completamente}. Nel caso della classificazione è il numero massimo di punti che si possono
classificare senza errori \emph{per tutte le possibili etichettature}.

Sia $X$ lo spazio delle istanze e si considerino $m$ punti. Attenzione: $m$ è il numero di punti considerati, che
può differire dal numero di dati $l$. Con due classi esistono $2^m$ possibili \termine{dicotomie}, cioè
etichettature dei punti con $-1$ o $+1$. Una dicotomia è \textbf{rappresentata} in $H$ se esiste $h\in H$ che la
realizza.

\begin{definizione}[Shattering]
$H$ \termine{frammenta} (\ingl{shatters}) un insieme di punti se e solo se può rappresentare \textbf{tutte} le
possibili dicotomie su di esso. Cioè, per ogni etichettatura esiste un'ipotesi $h\in H$ consistente con essa
(errore zero).
\end{definizione}

\begin{definizione}[VC-dimension]
La \termine{VC-dimension} di una classe di funzioni $H$ è la \textbf{cardinalità massima di un insieme di punti}
(in una qualche configurazione) che può essere frammentato da $H$. $\VC(H) = p$ significa che
\begin{itemize}
  \item $H$ frammenta \textbf{almeno un} insieme di $p$ punti;
  \item $H$ \textbf{non} frammenta \textbf{nessun} insieme di $p+1$ punti.
\end{itemize}
Se si possono frammentare insiemi arbitrariamente grandi, $\VC(H) = \infty$.
\end{definizione}

\begin{attenzione}[Asimmetria della definizione]
Per il limite inferiore basta mostrare \textbf{una} configurazione frammentabile. Per il limite superiore bisogna
mostrare che \textbf{nessuna} configurazione di $p+1$ punti lo è.
\end{attenzione}

\subsection{La VC-dimension delle rette nel piano}

\begin{teorema}[Rette nel piano]
Per gli iperpiani separatori (LTU) in $\R^2$, $h(\x) = \sign(\w\T\x+w_0)$, vale $\VC(H) = 3$.
\end{teorema}
\begin{proof}
\textbf{$\VC(H)\ge3$.} Si prendano 3 punti non allineati. Le loro $2^3=8$ dicotomie sono tutte separabili da una
retta (\figrif{fig:dicotomie}). Quando i tre punti hanno la stessa etichetta, basta una retta che li lasci tutti
dalla stessa parte. Altrimenti un punto ha etichetta diversa dagli altri due. Poiché i punti non sono allineati,
esiste una retta che isola quel vertice del triangolo dagli altri due. Non tutte le configurazioni di 3 punti
funzionano: tre punti allineati, con quello centrale di classe diversa, non sono separabili. Ma ne basta una.

\textbf{$\VC(H)<4$.} Dati 4 punti qualsiasi del piano, si distinguono tre casi, che coprono tutte le
possibilità.
\begin{itemize}
  \item \textbf{Tre dei punti sono allineati.} Si assegna al punto centrale dei tre una classe diversa da quella
        dei due estremi (\figrif{fig:separabilita}).
  \item \textbf{Nessuna terna è allineata e i 4 punti formano un quadrilatero convesso}: le due diagonali si
        incrociano. Si assegnano le estremità di una diagonale alla classe $+1$ e quelle dell'altra alla
        classe $-1$.
  \item \textbf{Nessuna terna è allineata e un punto è interno al triangolo formato dagli altri tre.} Si assegna
        al punto interno una classe e agli altri tre la classe opposta.
\end{itemize}
In ciascun caso l'etichettatura scelta non è separabile da una retta. Un semipiano è convesso: se contiene due
punti, contiene tutto il segmento che li unisce. Nel primo caso il semipiano che contiene i due estremi contiene
quindi anche il punto centrale. Nel secondo caso i segmenti che uniscono i punti della stessa classe si
incrociano (è uno XOR, \figrif{fig:quattro-punti}). Il punto d'incrocio dovrebbe stare in entrambi i semipiani,
che invece sono disgiunti. Nel terzo caso ogni semipiano che contiene i tre vertici contiene anche il triangolo, e
quindi il punto interno. Nessun insieme di 4 punti è frammentato.
\end{proof}

\begin{figure}[ht]
\centering
\begin{tikzpicture}[scale=0.62]
\foreach \a/\b/\c/\L [count=\i from 0] in {0/0/0/n, 1/0/0/A, 0/1/0/B, 0/0/1/C, 1/1/0/C, 1/0/1/B, 0/1/1/A, 1/1/1/n} {
  \pgfmathtruncatemacro{\col}{mod(\i,4)}
  \pgfmathtruncatemacro{\row}{\i/4}
  \begin{scope}[xshift=\col*3.5cm, yshift=-\row*3.1cm]
    \draw[grigio!60] (-0.6,-0.7) rectangle (2.6,2.2);
    \ifnum\a=1 \fill[blu] (0,0) circle (4.5pt); \else \draw[rosso, thick] (0,0) circle (4.5pt); \fi
    \ifnum\b=1 \fill[blu] (2,0) circle (4.5pt); \else \draw[rosso, thick] (2,0) circle (4.5pt); \fi
    \ifnum\c=1 \fill[blu] (1,1.6) circle (4.5pt); \else \draw[rosso, thick] (1,1.6) circle (4.5pt); \fi
    \if\L A \draw[thick] (0.85,-0.55) -- (-0.45,1.0); \fi
    \if\L B \draw[thick] (1.15,-0.55) -- (2.45,1.0); \fi
    \if\L C \draw[thick] (-0.45,0.95) -- (2.45,0.95); \fi
    \if\L n \draw[thick] (-0.45,-0.45) -- (2.45,-0.45); \fi
  \end{scope}
}
\end{tikzpicture}
\caption{Le $2^3=8$ dicotomie di tre punti non allineati (pieno $=+1$, vuoto $=-1$), ciascuna con una retta che
la realizza: le rette frammentano questi tre punti.}
\label{fig:dicotomie}
\end{figure}

\begin{figure}[ht]
\centering
\begin{tikzpicture}[scale=1.0]
\fill[blu] (0,0) circle (3pt); \fill[blu] (2,1.6) circle (3pt);
\draw[rosso, thick] (2,0) circle (3pt); \draw[rosso, thick] (0,1.6) circle (3pt);
\draw[blu] (0,0) -- (2,1.6); \draw[rosso] (2,0) -- (0,1.6);
\draw[densely dashed, grigio] plot[smooth] coordinates {(-0.4,0.9) (0.6,0.8) (1,0.4) (1.4,0.8) (2.4,0.7)};
\end{tikzpicture}
\caption{Con 4 punti esiste sempre un'etichettatura non separabile da una retta: qui i segmenti che uniscono i
punti della stessa classe si incrociano (XOR) e serve un confine non lineare.}
\label{fig:quattro-punti}
\end{figure}

\begin{teorema}[VC-dimension degli iperpiani]
La VC-dimension della classe degli iperpiani separatori (LTU) in uno spazio $n$-dimensionale è $n+1$.
\end{teorema}

È un valore massimo. Vincolando il modello può diminuire, ad esempio limitando la norma dei pesi come nelle SVM
(\rif{sec:svm-srm}).

\subsection{Altri esempi}

\begin{table}[ht]
\centering\small
\caption{VC-dimension di alcune classi di ipotesi.}
\begin{tabular}{lc}
\toprule
Classe $H$ & VC \\
\midrule
Soglie su $\R$, $\sign(w_1x+w_0)$ & 2 \\
Intervalli su $\R$ & 2 \\
Cerchi centrati nell'origine in 2D, $\sign(\x\T\x - b)$ & 1 \\
Cerchi centrati nell'origine, $\sign(k\,\x\T\x - b)$ & 2 \\
Cerchi in 2D (raggio e centro liberi) & 3 \\
Cerchi in 2D (raggio fisso, centro libero) & 3 \\
\bottomrule
\end{tabular}
\end{table}

Per i cerchi centrati nell'origine con il solo parametro $b$, il cerchio include i punti vicini all'origine. Con
due punti a distanze diverse non si può etichettare ``$+$'' il più lontano e ``$-$'' il più vicino, quindi
$\VC=1$. Aggiungendo $k$ si può scambiare l'interno con l'esterno, e la VC-dimension sale a 2.

\subsection{VC-dimension e numero di parametri}

La VC-dimension è legata al numero di parametri liberi, ma \textbf{non coincide} con esso:
\begin{itemize}
  \item nell'esempio dei cerchi, classi con lo \textbf{stesso} numero di parametri hanno VC-dimension diverse; e
        viceversa;
  \item si possono aggiungere parametri \textbf{ridondanti} senza cambiare nulla;
  \item esistono modelli con \textbf{un solo parametro e VC-dimension infinita}. L'esempio classico è
        $\sign(\sin(\alpha x))$: con un $\alpha$ abbastanza grande frammenta qualsiasi insieme di punti
        opportunamente scelti;
  \item per il \textbf{nearest neighbor} la VC-dimension è \textbf{infinita}: il 1-NN classifica correttamente
        qualunque etichettatura del training set.
\end{itemize}
La VC-dimension può essere difficile da calcolare per modelli con parametri non lineari.

\section{Il bound sul rischio}
\label{sec:vc-bound}

\begin{teorema}[VC-bound]
Con probabilità almeno $1-\delta$, per ogni $h\in H$ (con $\VC<l$, dove $l$ è il numero di dati):
\begin{equation}\label{eq:vc-bound}
\boxed{\;\underbrace{R[h]}_{\text{rischio}} \;\le\; \underbrace{\Remp[h] + \eps(\VC,l,\delta)}_{\text{rischio garantito}}\;}
\end{equation}
dove $\eps$ è la \termine{VC-confidence}. Ad esempio, per la loss 0/1:
\[
\eps(\VC,l,\delta) = \sqrt{\frac{\VC\left(\ln\frac{2l}{\VC}+1\right) - \ln\frac{\delta}{4}}{l}} .
\]
\end{teorema}

$\delta$ è la \textbf{confidenza}: con $\delta=0{,}01$ il bound vale con probabilità $0{,}99$. Esistono
formulazioni diverse per diverse classi di funzioni e task, scritte in generale come
$R\le\Remp+\eps\left(\frac1l,\VC,\frac1\delta\right)$ o $R\le\Remp+\eps\left(\frac{\VC}{l},-\frac{\ln\delta}{l}\right)$.
La VC-confidence \textbf{cresce con la VC-dimension} e \textbf{decresce con il numero di dati}. A meno di
logaritmi dipende dal rapporto $\VC/l$: raddoppiare i dati ha lo stesso effetto di dimezzare la complessità. Il
bound permette di \textbf{stimare l'errore sui dati futuri} usando solo l'errore di training e la VC-dimension. E
$\eps$ si può calcolare \textbf{prima} dell'apprendimento.

\begin{idea}[Il bound spiega underfitting e overfitting]
\begin{itemize}
  \item \textbf{Più dati} ($l$ alto): $\eps$ diminuisce e il bound si avvicina a $R$. Il rischio empirico diventa
        una buona stima di quello vero.
  \item \textbf{Modello troppo semplice} (VC bassa): $\eps$ piccolo ma $\Remp$ alto: \textbf{underfitting}.
  \item \textbf{Modello troppo complesso} (VC alta, $l$ fissato): $\Remp$ basso ma $\eps$ alto, e con esso il
        bound su $R$: \textbf{overfitting}.
\end{itemize}
In parole semplici, si può approssimare bene $f$ dagli esempi \uline{purché si abbiano abbastanza dati e la
complessità del modello sia adatta al problema}. Bisogna adattarsi ai dati abbastanza da evitare l'underfitting.
Ma non troppo, per evitare l'overfitting.
\end{idea}

\begin{figure}[ht]
\centering
\begin{tikzpicture}
\begin{axis}[stile, width=8.6cm, height=5.4cm, domain=0.6:10, xmin=0, xmax=10.6, ymin=0, ymax=1.15,
  xtick=\empty, ytick=\empty, xlabel={VC-dimension (complessità di $H$)}, ylabel={errore}, legend pos=north east,
  legend style={at={(0.98,0.98)}}]
\addplot[thick, blu] {0.95*exp(-x/2.2)}; \addlegendentry{$\Remp$ (training)}
\addplot[thick, rosso] {0.035*x^1.35}; \addlegendentry{VC-confidence $\eps$}
\addplot[very thick, viola] {0.95*exp(-x/2.2)+0.035*x^1.35}; \addlegendentry{bound su $R$}
\draw[densely dashed, grigio] (axis cs:3.82,0) -- (axis cs:3.82,0.34);
\node[font=\footnotesize, above, viola] at (axis cs:3.82,0.42) {minimo del bound};
\node[font=\footnotesize, grigio] at (axis cs:1.3,0.15) {underfitting};
\node[font=\footnotesize, grigio] at (axis cs:8.6,0.25) {overfitting};
\end{axis}
\end{tikzpicture}
\caption{A numero di dati fissato: il rischio empirico decresce con la VC-dimension, la VC-confidence cresce; la
loro somma, il bound su $R$, ha un minimo nel punto di miglior compromesso (andamento qualitativo).}
\label{fig:vc-bound}
\end{figure}

\begin{esempio}[Ritrovare $h$ e $h'$]
Nella definizione di overfitting (\rif{sec:overfitting}), $h$ ha errore di training minore ma rischio maggiore di
$h'$. Sul grafico della \figrif{fig:vc-bound}, $h$ sta \textbf{a destra del minimo}: VC più alta, $\Remp$ più
basso, bound più alto. $h'$ sta vicino al minimo.
\end{esempio}

\begin{attenzione}[Bound ``nel caso peggiore'']
I bound devono valere per tutte le funzioni di $H$ e tutti i training set, tranne una frazione $\delta$. Sono
quindi \ingl{worst case} e spesso molto larghi. La VC-confidence può essere \textbf{centinaia di volte} più grande
dell'effetto di overfitting osservato empiricamente.
\end{attenzione}

\begin{idea}[Una conseguenza pratica]
Per molte classi ragionevoli (ad esempio gli approssimatori lineari) la VC-dimension è \textbf{lineare nel numero
di parametri liberi}. Quindi per apprendere ``bene'' serve un numero di esempi \textbf{lineare nella
VC-dimension}, cioè nel numero di parametri. È una giustificazione teorica di euristiche come ``servono almeno
$k$ esempi per parametro''.
\end{idea}

\section{Structural Risk Minimization}
\label{sec:srm}

\begin{definizione}[Structural Risk Minimization]
La \termine{SRM} usa la VC-dimension come parametro di controllo. Invece di minimizzare solo $\Remp$,
\textbf{minimizza il bound}, cioè la somma di errore di training e VC-confidence. Si definisce una
\textbf{struttura annidata} di spazi delle ipotesi ordinati per VC-dimension (finita)
\[
H_1\subseteq H_2\subseteq\dots\subseteq H_n,\qquad \VC(H_1)\le\dots\le\VC(H_n),
\]
e si sceglie lo spazio (e in esso l'ipotesi) con il \textbf{miglior bound sul rischio vero}.
\end{definizione}

La SRM formalizza il \textbf{controllo della complessità}, cioè un compromesso tra complessità (VC-dimension) e
accuratezza sul training (\ingl{fitting}). Esempi di strutture annidate:
\begin{itemize}
  \item reti neurali con numero crescente di \textbf{unità nascoste} (circa). Ma anche con numero crescente di
        \textbf{epoche}, perché conta anche la grandezza dei pesi (\rif{sec:vc-reti});
  \item polinomi di \textbf{grado} $M$ crescente;
  \item valori crescenti di $c$ con il vincolo $\norm{\w}<c$, cioè la \textbf{regolarizzazione}. È controllata
        equivalentemente da $\lambda$: aumentare $\lambda$ equivale a spostarsi verso spazi più piccoli;
  \item alberi di decisione con numero crescente di nodi (circa).
\end{itemize}
Gli \textbf{iperparametri importanti} di un modello sono proprio quelli legati alla VC-dimension.

\subsection{A cosa serve il bound}

\begin{itemize}
  \item Dà un \textbf{fondamento teorico} al ML, indipendente dai dettagli del modello e dell'algoritmo. Il
        rischio si può limitare analiticamente, e il controllo della complessità ha un ruolo preciso. La scelta
        ottima della complessità sulla struttura minimizza il rischio atteso: è il \textbf{principio induttivo
        della SRM}. Il bound spiega anche la differenza tra il ML e l'ottimizzazione pura: nel ML l'obiettivo non
        è minimizzare l'errore sui dati, ma sui dati \emph{futuri}.
  \item Come \textbf{stima dell'errore} è usato raramente:
        \begin{itemize}
          \item può essere molto largo. Resta però utile per la \textbf{model selection}, dove conta la
                \textbf{dimensione relativa} (il \emph{ranking}) degli errori;
          \item è un upper bound probabilistico valido per tutte le funzioni della classe; quindi si può cercare
                nella classe;
          \item la VC-dimension di classi specifiche è difficile da calcolare;
          \item il bound grezzo può non essere adeguato per una valutazione affidabile (\textbf{model
                assessment}). Si studiano bound più stretti.
        \end{itemize}
  \item Indica una direzione per sviluppare \textbf{nuovi modelli} guidati dalla SRM.
\end{itemize}

\subsection{Due approcci pratici alla SRM}

\begin{enumerate}
  \item \textbf{Fissare la complessità} (quindi la VC-confidence) e \textbf{minimizzare l'errore di training}.
        È l'approccio delle reti neurali. Alcune euristiche di addestramento realizzano già una SRM
        \textbf{implicita} (l'early stopping). Oppure la si realizza con la \textbf{regolarizzazione di
        Tikhonov}, minimizzando $\Remp$ più un termine di complessità.
  \item \textbf{Fissare l'errore di training} e \textbf{minimizzare automaticamente la VC-confidence}. È l'idea
        delle \textbf{SVM} (capitolo~\ref{cap:svm}).
\end{enumerate}

\begin{attenzione}[Oltre il compromesso classico]
Fenomeni più recenti, tipici delle reti profonde sovra-parametrizzate (\ingl{double descent},
\ingl{benign overfitting}), sembrano sfidare questa visione. Sono discussi nel \rif{sec:double-descent}.
\end{attenzione}

\begin{puntichiave}
\begin{itemize}
  \item ERM: si minimizza $\Remp$ al posto di $R$; la domanda è quanto $\Remp$ approssimi $R$.
  \item VC-dimension: il massimo numero di punti che $H$ frammenta; $n+1$ per gli iperpiani in $\R^n$; non
        coincide con il numero di parametri (infinita con un solo parametro, infinita per il NN).
  \item VC-bound: $R\le\Remp+\eps(\VC,l,\delta)$, con $\eps$ crescente in VC e decrescente in $l$.
  \item SRM: su una struttura annidata $H_1\subseteq H_2\subseteq\dots$ si sceglie lo spazio con il bound
        minimo. Due realizzazioni: complessità fissa + minimo $\Remp$ (Tikhonov), $\Remp$ fisso + minima
        VC-confidence (SVM).
\end{itemize}
\end{puntichiave}

\begin{domande}
\begin{enumerate}
  \item Differenza tra rischio $R$ e rischio empirico $\Remp$; cos'è il principio ERM?
  \item Definire shattering e VC-dimension; dimostrare che la VC-dimension delle rette nel piano è 3.
  \item La VC-dimension coincide con il numero di parametri liberi? Esempi.
  \item Scrivere il VC-bound e spiegare come giustifica underfitting e overfitting.
  \item Cos'è la Structural Risk Minimization? Esempi di strutture annidate. Limiti pratici del bound.
\end{enumerate}
\end{domande}
